% frontmatter/polish/como-usar.tex -- "Jak korzystać z tego zeszytu":
% las instrucciones para el adulto de frontmatter/como-usar.tex, en
% polaco.
\section*{Jak korzystać z tego zeszytu}

Ten zeszyt jest dla dziecka w wieku pięciu albo sześciu lat, które
zaczyna poznawać liczby: rozpoznawać je, liczyć i pisać -- a z czasem
także porównywać, dodawać i odejmować. Należy do tej samej rodziny co
hiszpańskie zeszyty \emph{Aprendo a leer}, z tymi samymi bohaterami
(Lucía, jej młodszy brat Dani, ich pies Toby, mama, tata i babcia Rosa).
A strona po stronie jest to polskie wydanie \emph{Aprendo los números}:
ta sama liczba, te same rysunki i to samo zadanie każdego dnia, więc
dziecko może korzystać z jednego zeszytu albo z obu.

Na każdy dzień roboczy w roku, od poniedziałku do piątku, także w
wakacje, jest jedna strona. Pięć albo dziesięć minut dziennie w
zupełności wystarczy. Każda strona wygląda tak samo:

\begin{itemize}[leftmargin=6mm]
\item \textbf{Data}, \textbf{Zrobione} i \textbf{Uwagi}: dla dorosłego.
  Zaznaczenie codziennie, czy zadanie było zrobione samodzielnie, z
  pomocą czy z trudem, zajmuje chwilę, a po roku jest zapisem
  wszystkiego, czego dziecko się nauczyło.
\item \textbf{Liczba dnia}, w zielonej ramce: duża liczba, jej nazwa,
  \textbf{ramka dziesiątkowa} -- dwa rzędy po pięć kratek, z kropką za
  każdą rzecz, żeby od razu było widać, ile to jest (5 to cały rząd; 7
  to rząd i jeszcze dwie) -- i tyle samo narysowanych rzeczy. Pod spodem
  zdanie z historyjki, które czyta dorosły: liczba dnia, w domu Lucíi.
\item \textbf{Zadanie}, które prawie zawsze kończy się kolorowaniem,
  zakreślaniem albo rysowaniem. Polecenia (małą, szarą czcionką)
  \textbf{czyta dorosły}.
\end{itemize}

Każde zadanie jest opisane, jedno po drugim, na następnej stronie.

\subsection*{Jak pomagać}

{\small
\begin{itemize}[leftmargin=6mm, itemsep=1pt, topsep=2pt]
\item Liczyć to dotykać: przy liczeniu wskazujemy palcem każdą rzecz,
  i tylko raz. Jeśli któraś zostanie pominięta albo policzona dwa razy,
  liczymy razem, powoli.
\item Mówimy nazwę liczby za każdym razem, kiedy dziecko ją pisze albo
  znajduje: liczbę się widzi, mówi i pisze.
\item Po polsku liczba zmienia się razem z tym, co liczymy (jeden pies,
  dwa psy, pięć psów; dwie piłki, troje dzieci). Przy liczeniu wystarczy
  mówić same liczby: jeden, dwa, trzy...
\item Jeśli coś sprawia trudność, można to powtórzyć następnego dnia z
  rzeczami z domu: łyżkami, guzikami, butami.
\item Na koniec każdej części jest medal, a na końcu zeszytu dyplom i
  \textbf{odpowiedzi} do stron, które je mają.
\end{itemize}}
\newpage
